\documentclass[11pt,a4paper]{article}

\usepackage[T1]{fontenc}
\usepackage[utf8]{inputenc}
\usepackage[slovene]{babel}
\usepackage{lmodern}
\usepackage{microtype}
\usepackage{geometry}
\usepackage{amsmath,amssymb}
\usepackage{enumitem}
\usepackage{fancyhdr}

\geometry{margin=1.8cm}
\setlength{\parindent}{0pt}
\setlength{\parskip}{0pt}
\setlist[enumerate]{leftmargin=*,itemsep=0pt,topsep=0pt,parsep=0pt}
\pagestyle{fancy}
\fancyhf{}
\lhead{\small Matematika}
\rhead{\small Preizkus znanja: kombinatorika}
\cfoot{\thepage}

\newenvironment{taskbox}[1]{%
    \par\addvspace{8pt}%
    \noindent\textbf{\large #1}\par
}{\par\addvspace{4pt}}
\newcommand{\work}[1]{\par\vspace{#1}}

\begin{document}

\begin{center}
    {\LARGE\textbf{Preizkus znanja: Kombinatorika}}\\[4pt]
    {\large Čas pisanja: 60 minut \qquad Skupaj: 40 točk}
\end{center}

\textbf{Ime in priimek:} \rule{8cm}{0.4pt}
\hfill
\textbf{Datum:} \rule{3cm}{0.4pt}

\smallskip
\textbf{Navodila:} Pri vsaki nalogi zapiši postopek in račun.
Rezultat brez ustreznega postopka lahko prinese manj točk.
Pri nalogah si lahko pomagate s kalkulatorjem. 

\begin{taskbox}{1. \hfill 10 točk}
\begin{enumerate}[label=\alph*)]
    \item Na plavalnem tekmovanju je 11 tekmovalcev. Na koliko načinov se lahko razdelijo
    zlata, srebrna in bronasta mesta? Vsak tekmovalec lahko zasede največ eno mesto.
    \hfill (4 t)
    \work{7cm}
    \item Sestavljamo petmestno številčno geslo (PIN) iz števk
    \(0,1,\ldots,9\). Števke se lahko ponavljajo in začetna števka je lahko 0.
    Koliko različnih PIN-kod je mogoče sestaviti? Utemelji račun. \hfill (6 t)
    \work{7cm}
\end{enumerate}
\end{taskbox}

\clearpage
\begin{taskbox}{2. \hfill 5 točk}
\begin{enumerate}[label=\alph*)]
    \item V slaščičarni lahko izberemo eno od 4 sadnih ali eno od 3 čokoladnih
    sladic. Koliko različnih sladic lahko izberemo? \hfill (2 t)
    \work{7cm}
    \item Za malico izberemo enega od 4 sendvičev in eno od 3 pijač. Nariši
    kombinatorično drevo ter zapiši število različnih malic. \hfill (3 t)
    \work{7cm}
\end{enumerate}
\end{taskbox}

\clearpage
\begin{taskbox}{3. \hfill 7 točk}
\begin{enumerate}[label=\alph*)]
    \item Iz črk besede \textsc{triglav} sestavljamo nove besede. Vsakič
    uporabimo vse črke natanko enkrat. Koliko je takih besed, v katerih vsi
    soglasniki stojijo skupaj?     \hfill (3 t)
    \work{10cm}
    \item Koliko različnih nizov lahko sestavimo s preurejanjem vseh črk besede
    \textsc{kaktus}? \hfill (4 t)
    \work{7cm}
\end{enumerate}
\end{taskbox}

\clearpage
\begin{taskbox}{4. \hfill 9 točk}
\begin{enumerate}[label=\alph*)]
    \item Razvij \((2x+3)^4\) v vsoto enočlenikov. \hfill (5 t)
    \work{7cm}
    \item Z uporabo formule za \((k+1)\)-ti člen določi tretji člen razvoja
    \((x-2)^8\). Pokaži, kako določiš \(k\). \hfill (4 t)
    \work{7cm}
\end{enumerate}
\end{taskbox}

\clearpage
\begin{taskbox}{5. \hfill 9 točk}
V skupini je 9 fantov in 7 deklet.
\begin{enumerate}[label=\alph*)]
    \item Izberemo petčlansko ekipo, v kateri sta natanko 2 dekleti. Na koliko
    načinov jo lahko izberemo? \hfill (5 t)
    \work{7cm}
    \item Izberemo štiričlansko skupino, v kateri je vsaj eno dekle. Na koliko
    načinov jo lahko izberemo? \hfill (4 t)
    \work{2.5cm}
\end{enumerate}
\end{taskbox}

\vfill
\begin{center}
    \textbf{Preden oddaš:} preveri, ali si pri vsaki nalogi upošteval vrstni red
    in možnost ponavljanja.
\end{center}

\end{document}
